\documentclass[11pt]{article}

\usepackage{amsmath, amssymb, amsthm, amsfonts, mathtools}
\usepackage[hidelinks]{hyperref}
\usepackage{xcolor}
\usepackage{geometry}
\geometry{a4paper, margin=1in}

\theoremstyle{plain}
\newtheorem{theorem}{Theorem}

\theoremstyle{remark}
\newtheorem{remark}{Remark}
\newtheorem{recommendation}{Recommendation}

\title{Novelty Assessment and Next Steps\\
\large For the Proof of Conjecture~3: The Period Length at $\omega = 1$}
\author{Dirk Kunert \thanks{Independent Researcher, Email: dirk.kunert@gmail.com}}
\date{\today}

\begin{document}

\maketitle

\begin{abstract}
We assess the novelty of the theorem $\lambda_{\alpha,\beta} = \alpha + \beta + 1$ for
one-dimensional cut-and-project sequences at window width $\omega = 1$.
An extensive survey of related literature --- including the three-distance theorem,
Christoffel word theory, and the cut-and-project literature on model sets ---
reveals no prior result that states or proves this formula.
We conclude that the theorem is almost certainly novel, identify the closest related
known results, and provide recommendations for publication.
\end{abstract}

% =====================
\section{The Theorem Under Review}
% =====================

The result whose novelty is assessed here is the following.

\begin{theorem}
Let $\alpha, \beta \in \mathbb{N}$ with $\gcd(\alpha, \beta) = 1$, and let $\omega = 1$.
Consider the lattice points $(x, y) \in \mathbb{Z}^2$ satisfying
\[
    \frac{\alpha}{\beta}(x - 1) \;\le\; y \;\le\; \frac{\alpha}{\beta}\,x + 1,
\]
projected vertically onto $f(x) = \tfrac{\alpha}{\beta}\,x$.
The period length of the resulting distance sequence $\bigl(\delta^{(i)}\bigr)$ equals
\[
    \lambda_{\alpha,\beta} \;=\; \Lambda_{\alpha,\beta} \;:=\; \alpha + \beta + 1.
\]
\end{theorem}

The proof (see \cite{Kunert2025proof}) proceeds in four steps:
\begin{enumerate}
    \item The invariant $r := \alpha x - \beta y$ takes values in
          $\mathcal{R} = \{-\beta, \ldots, 0, \ldots, \alpha\}$,
          giving $|\mathcal{R}| = \alpha + \beta + 1$ distinct values.
    \item For each fixed $r \in \mathcal{R}$, the projected values
          $\tilde{p} = \beta x + \alpha y$ form an arithmetic progression with
          common difference $\alpha^2 + \beta^2$.
    \item By B\'{e}zout's identity, the $\alpha + \beta + 1$ residues modulo
          $\alpha^2 + \beta^2$ are pairwise distinct.
    \item The sorted union of these progressions therefore has period exactly
          $\alpha + \beta + 1$.
\end{enumerate}

% =====================
\section{Related Literature}
% =====================

\subsection{The Three-Distance Theorem}

The three-gap theorem (Steinhaus, S\'{o}s, Sur\'{a}nyi, \'{S}wierczkowski; 1950s) states
that placing $n$ points $\{k\alpha\}$ on a unit circle produces at most three distinct
gap lengths \cite{ThreeGap}.
When $\alpha = p/q$ is rational the number of distinct gap lengths reduces to at most two.
More recent treatments include the lattice proof by Marklof and Str\"{o}mbergsson (2017)
and the exposition by Berth\'{e} and Reutenauer \cite{BertheReutenauer2024}.

This theorem is related in spirit but does not imply the result above for two reasons:
\begin{itemize}
    \item It concerns rotations on a circle, not lattice projections through a
          strip of finite width.
    \item It characterises the \emph{number} of distinct gap lengths (at most three),
          not the \emph{period length} of the gap sequence.
          The quantity $\alpha + \beta + 1$ does not appear in this framework.
\end{itemize}

\subsection{Christoffel Words and Sturmian Sequences}

The closest structural relative found in the literature is the theory of Christoffel
words \cite{BerstelEtAl2008}.
The lower Christoffel word of slope $p/q$ with $\gcd(p, q) = 1$ is a word of length
$p + q$ over a two-letter alphabet, and the mechanical sequence of rational slope $p/q$
is periodic with period $\mathbf{p + q}$.

The difference from Theorem~1 is mathematically significant:

\begin{itemize}
    \item Christoffel words have period $p + q$; the theorem gives period
          $\alpha + \beta + \mathbf{1}$.
          The extra $+1$ is a direct consequence of the symmetric window of width
          $\omega = 1$, which allows $r$ to range over $\alpha + \beta + 1$ values
          instead of $\alpha + \beta$.
    \item The Christoffel/Sturmian framework describes a \emph{single} line
          (window of vanishing width), not a strip of finite width $\omega = 1$
          capturing lattice points on both sides of the central line.
    \item The object of study here is the sequence of \emph{Euclidean distances}
          between projected points, not a symbolic word period.
\end{itemize}

\subsection{Cut-and-Project Literature (Model Sets)}

The mainstream literature on model sets and quasicrystals
\cite{Senechal2009, HaynesEtAl2016} focuses almost exclusively on irrational slopes,
since rational slopes produce trivially periodic sequences.
Key topics in this literature --- linear repetitivity, patch frequencies, diffraction
spectra --- differ from the question addressed here.
No paper was found that determines the period of the distance sequence for a
one-dimensional cut-and-project set with rational slope and a symmetric window of
finite width.

\subsection{Beatty Sequences and Fraenkel's Conjecture}

The Beatty sequence and Fraenkel conjecture literature concerns partitions of
$\mathbb{N}$ into Beatty sequences with rational or irrational parameters.
None of the papers in this area address projection distances with a symmetric window
of the kind appearing in Theorem~1.

% =====================
\section{Assessment of Novelty}
% =====================

Based on the survey above, the theorem is \textbf{almost certainly novel}.
The specific reasons are:

\begin{enumerate}
    \item \textbf{The formula $\alpha + \beta + 1$ does not appear in the existing
          literature} in this context.
          The nearest analogue --- the period $p + q$ of Christoffel words --- differs
          by exactly $1$, and this difference is mathematically consequential.

    \item \textbf{The construction is not a standard object of study.}
          Most cut-and-project work concerns irrational slopes or asymptotic behaviour.
          The rational slope case with a symmetric window of width $\omega = 1$ is not
          treated in the mainstream literature.

    \item \textbf{The proof strategy is original.}
          Classifying the $\alpha + \beta + 1$ arithmetic progressions in the projected
          coordinate $\tilde{p} = \beta x + \alpha y$ via the invariant
          $r = \alpha x - \beta y$, and then invoking B\'{e}zout's identity to establish
          distinct residues modulo $\alpha^2 + \beta^2$, does not appear to have been
          written down before.
\end{enumerate}

\begin{remark}
    A full MathSciNet/zbMATH/Scopus search by a librarian with database access would
    be a prudent final step before submission.
    Web-based searches cannot access all specialised journal content.
\end{remark}

% =====================
\section{Recommendations for Publication}
% =====================

\begin{recommendation}
Add a paragraph to the introduction of \cite{Kunert2025proof} explaining why
Theorem~1 is \emph{not} a direct consequence of Christoffel word theory.
The key points to make are:
\begin{itemize}
    \item The period of a Christoffel word of slope $\alpha/\beta$ is $\alpha + \beta$,
          not $\alpha + \beta + 1$.
    \item The extra $+1$ arises from the symmetric window $\omega = 1$, which widens
          the strip and adds one additional residue class.
    \item The result concerns Euclidean distances in the projected sequence, not a
          symbolic word over a finite alphabet.
\end{itemize}
\end{recommendation}

\begin{recommendation}
Cite the following works as the closest known results, to help a reviewer see the gap
between the existing literature and the new theorem:
\begin{itemize}
    \item Berth\'{e} and Reutenauer \cite{BertheReutenauer2024} on the three-distance
          theorem.
    \item Berstel, Lauve, Reutenauer, and Saliola \cite{BerstelEtAl2008} on
          Christoffel words.
\end{itemize}
\end{recommendation}

\begin{recommendation}
Request a zbMATH/MathSciNet search using the terms
``cut-and-project'', ``period length'', ``distance sequence'',
``rational slope'', and ``coprime''
before final submission, to rule out any overlooked specialised result.
\end{recommendation}

% =========================
\begin{thebibliography}{9}
% =========================

\bibitem{Kunert2025proof}
Dirk Kunert, 2025,
\emph{Proof of Conjecture~3: The Period Length at $\omega = 1$},
available at: \url{https://github.com/dkunert/cut-and-project}.

\bibitem{ThreeGap}
\emph{Three-gap theorem},
Wikipedia,
available at: \url{https://en.wikipedia.org/wiki/Three-gap_theorem},
accessed April 2025.

\bibitem{BertheReutenauer2024}
Val\'{e}rie Berth\'{e} and Christophe Reutenauer,
\emph{On the Three Distance Theorem},
\emph{The Mathematical Intelligencer}, 2024.
Available at: \url{https://link.springer.com/article/10.1007/s00283-023-10316-z}.

\bibitem{BerstelEtAl2008}
Jean Berstel, Aaron Lauve, Christophe Reutenauer, and Franco V.\ Saliola,
\emph{Combinatorics on Words: Christoffel Words and Repetitions in Words},
CRM Monograph Series, Vol.~27,
American Mathematical Society, 2008.

\bibitem{Senechal2009}
Marjorie Senechal,
\emph{Quasicrystals and Geometry},
Cambridge University Press, 2009.
ISBN: 978-0-521-57541-6.

\bibitem{HaynesEtAl2016}
Alan Haynes, Henna Koivusalo, James Walton, and Lorenzo Sadun,
\emph{Gaps Problems and Frequencies of Patches in Cut and Project Sets},
\emph{Mathematical Proceedings of the Cambridge Philosophical Society}, 2016.
Available at: \url{https://www.cambridge.org/core/journals/mathematical-proceedings-of-the-cambridge-philosophical-society/article/abs/gaps-problems-and-frequencies-of-patches-in-cut-and-project-sets/1A47C2BE5E660C5C0A0AC4A828F3B6E7}.

\end{thebibliography}

\end{document}
